\section{Backlog and future work}\label{sec:backlog}
What is deliberately left undone, and what has to exist before each item can be written down.
\begin{itemize}
    \item \textbf{Commanded gimbal control.} A time-varying $\theta_t(t)$ driven by a control law;
    a 3D item, since steering needs an out-of-plane degree of freedom. \emph{Waiting on:} a
    guidance objective, actuator limits, a sensor model, and the 3D model. The air brake is a
    different loop, implemented in 2D with a proportional law; a law with integral or derivative
    memory on the zero-order-hold mode (signal, gains, saturation, anti-windup) is unwritten, as
    are brake retraction and a rate limit.
    \item \textbf{Wing-body aerodynamics.} Replace the constant, linearized and table
    coefficient models with a Barrowman-style build-up: each of nose, body and fins contributes a
    normal-force slope $C_{z,\alpha,i}$ at a location $X_i$, and the centre of pressure is the
    weighted average $X_{cp}=\sum_iC_{z,\alpha,i}X_i/\sum_iC_{z,\alpha,i}$, so that $l_p-l_g$ is
    computed rather than read from a table. A Mach correction is needed beyond low subsonic speed.
    \emph{Waiting on:} an actual airframe --- nose shape, body diameter and length, fin planform.
    \item \textbf{Dynamic stability, quantitatively.} The damped short-period framework exists;
    the value of $C_{m,q}$ does not. \emph{Waiting on:} a measured value (forced-oscillation
    tests or unsteady CFD) or a distributed-force derivation from the airframe.
    \item \textbf{Propellant grain geometry} behind $I_{m,g}(t)$ and $l_{g,m}(t)$: choose a grain
    family and its dimensions, from which $A_b(w)$, $l_{g,m}(t)$ and $I_{m,g}(t)$ follow by
    integration.
    \item \textbf{Component-based mass properties} of the dry rocket, built from nose cone, tubes,
    fins and internal parts with their densities. \emph{Waiting on:} the airframe geometry and a
    materials table.
    \item \textbf{Vertical wind.} Terrain-slope orographic lift or a convective parameterisation;
    each needs a data source the project does not have. On hold pending research.
    \item \textbf{CFD pipeline} for $C_x,C_z,C_m$: the run matrix or wind-tunnel campaign that
    populates the aerodynamic table for a specific airframe.
    \item \textbf{Out-of-plane motion}: crosswind weathercocking and the lateral Coriolis
    deflection, both lifted by a full 6DOF model.
    \item \textbf{State-coupled turbulence.} A Dryden filter driven by the true airspeed, carrying
    its own state on the sampled mode. \emph{Waiting on:} its discretisation, and whether per-tick
    noise makes results depend on the sample rate.
    \item \textbf{Rail--rocket dynamics, in detail.} Derive the reaction (normal) forces at the
    rail buttons and the friction they cause, and the tip-off transient at rail exit. \emph{Waiting
    on:} button positions, clearances and a friction coefficient.
    \item \textbf{Rocket and parachute dynamics, in detail.} A coupled two-body model: pendulum
    motion (wobble) of the rocket under the canopy, riser elasticity, and the canopy opening
    transient with its inflation shock. \emph{Waiting on:} canopy and riser data, and the
    attitude degree of freedom under canopy.
    \item \textbf{Roll dynamics with a reaction wheel (2D).} A decoupled roll degree of freedom
    about $\hat{x}_r$: roll angle $\gamma_r$, roll rate $\omega_{x}$ and axial moment of inertia
    $I_x$. A reaction wheel torques the airframe with an equal and opposite torque on the wheel,
    whose speed and angular momentum saturate. Roll is also damped and driven aerodynamically
    (roll damping, roll moment from fin cant).
    \emph{Waiting on:} $I_x$ of the airframe, the wheel inertia and motor torque limits, the
    aerodynamic roll coefficients, and a control law.
\end{itemize}
